\PassOptionsToPackage{dvispnames, table}{xcolor}
\documentclass[aspectratio=169, hyperref={colorlinks=true}]{beamer}
\usepackage[beamer]{prettytex/base}
\usepackage{prettytex/math}
\usepackage{prettytex/gfx}
\usepackage[backend=biber, citestyle=numeric, bibstyle=numeric, hyperref=true, sorting=none, maxbibnames=99]{biblatex}
\usepackage[toc, acronym, style=long3col, indexonlyfirst=true, nogroupskip=true]{glossaries}
\usepackage{csquotes}
\usepackage{fontawesome}
\addbibresource{../literature/sources.bib}
\renewcommand{\figureloc}{../figures}  % for tikzfigures
\graphicspath{{../figures}}  % for all other figures
\usetheme{oxonian}

\title{General Kernel Spectral Methods for Equilibrium Measures \\ \normalsize MMSC Dissertation}
\titlegraphic{\includegraphics[width=2cm]{../figures/logos/ox-circle.eps}}
\author{Peter Julius Waldert}
\institute{Mathematical Institute \\ University of Oxford}
\date{$18^{\rm th}$ of September, 2023}

\definecolor{lighterOxfordBlue}{RGB}{39, 63, 111}
\definecolor{oxfordGolden}{RGB}{175, 148, 72}
\definecolor{wongGreen}{RGB}{101, 145, 87}
\definecolor{wongPurple}{RGB}{204, 121, 167}
\definecolor{powderBlue}{RGB}{158, 179, 194}
\setbeamercolor{block title}{fg=white, bg=lighterOxfordBlue} % or head
\setbeamercolor{block body}{parent=normal text, bg=lightgrey} % or body, foot
\setbeamercolor{alerted text}{fg=wongGreen}
\renewcommand{\operatorcolor}{lighterOxfordBlue}

% Also don't forget your slides for the first part of the viva when the examiners will ask you to describe in about 5 minutes the main aims and results of your dissertation. They recommend the slides might summarise: the problem formulation (the problem you are trying to solve and its mathematical formulation); analysis and results; any original contributions you have made; challenges faced and future work.

\makenoidxglossaries
\newacronym{fcc}{fcc}{face-centered cubic}
\newacronym{ml}{ML}{Machine Learning}
\newacronym{gd}{GD}{Gradient Descent}
\newacronym{gui}{GUI}{Graphical User Interface}
\newacronym{lru}{LRU}{Least Recently Used}
\newacronym{pde}{PDE}{Partial Differential Equation}
\newacronym{ad}{AD}{Automatic Differentiation}
\newacronym{lbfgs}{LBFGS}{Limited-memory Broyden-Fletcher-Goldfarb-Shanno}

\usepackage{xifthen}
\newcommand{\name}[1]{\textsc{#1}}
\newcommand{\inv}{^{-1}}
\newcommand{\qRq}{\quad\Rightarrow\quad}
\newcommand{\qLRq}{\quad\Leftrightarrow\quad}
\newcommand{\functionspace}{L}
\newcommand{\functionspacehat}{\hat{L}}
\renewcommand{\norm}[1]{\left\lVert#1\right\rVert_{\scriptscriptstyle 2}}
\newcommand{\jacobiarg}[1]{2 \norm{\vec{#1}}^2 - 1}
\newcommand{\jacobi}[2][]{P_{\ifthenelse{\isempty{#1}}{k}{#1}}^{(a, b)}\left(\jacobiarg{#2}\right)}
\newcommand{\jacobivec}[1]{\vec{P}^{(a, b)}\left(\jacobiarg{#1}\right)}
\newcommand{\weight}[2][]{\left(1-\norm{\vec{#2}}^2\right)^{\ifthenelse{\isempty{#1}}{a}{#1}}}
\newcommand{\goesto}{\rightarrow}
\newcommand{\hatvec}[1]{\hat{\vec{#1}}}

\newcommand{\hierKoennteIhreWerbungStehen}{
  \begin{center}
    \bfseries \textcolor{themecolor2}{To be done / included.}
  \end{center}
}


\begin{document}
  {\setbeamertemplate{footline}{}\frame{\titlepage}}

  \begin{frame}{Connection to Swarming Behaviour}
    \begin{figure}[H]
      \centering
      \includegraphics[width=0.56\linewidth]{results/morse-2d/simulation-quiver.pdf}
      \caption[Quiver plot of 120 particles in 2D interacting through the Morse potential]{$N_p = 120$ particles in $d = 2$ dimensions. Morse potential, friction and self-propulsion parameters are as given in \cite{2006-self-propelled}, reproducing their results.}
      \label{fig:simulation-quiver-illustration}
    \end{figure}
  \end{frame}

  \begin{frame}{Problem Setting}
    \begin{definition}[Equilibrium Measure]
      For a given pairwise interaction potential $K: \R \to \R$, the equilibrium measure $\hat{\rho}: D \to \R$ with $D \subseteq \R^d$ is a measure chosen such that
      $$U_K[\hat{\rho}] := \frac{1}{2} \iint K\left(\norm{\hatvec{x} - \hatvec{y}}\right) \,\dd\hat{\rho}(\hatvec{x})\,\dd\hat{\rho}(\hatvec{y})\,,$$
      is minimised, where $\dd\hat{\rho} = \hat{\rho}(\hatvec{x})\dd\hatvec{x}$ \parencite{2021-arbitrary-dimensions}.
    \end{definition}
    \pause

    For example: in a two-particle system, $\hat{\rho}(\hatvec{x}) = \delta(\hatvec{x}-\hatvec{p_1}) + \delta(\hatvec{x}-\hatvec{p_2})$ to see that
    $U_K[\hat{\rho}] = K(0) + K(\norm{\hatvec{p_1} - \hatvec{p_2}})\,,$
    and hence, the total energy becomes
    $$E = \frac{m}{2} \left(\norm{\hatvec{v_1}}^2 + \norm{\hatvec{v_2}}^2\right) + K(0) + K(\norm{\hatvec{p_1} - \hatvec{p_2}})\,.$$
  \end{frame}

  \begin{frame}{Problem Formulation}
    \begin{definition}[Particle Density Distribution Problem]
      Given an interaction kernel $K: \R^+ \to \R$, the density distribution problem is to find the equilibrium measure $\hat{\rho}: B_R(\vec{0}) \to \R$ of mass $M = 1$ on a $d$-dimensional ball of radius $R \in \R^+$ that minimises the total potential $U_K[\hat{\rho}]$.
    \end{definition}
    \pause

    We then consider a \alert{spectral method} with the following ansatz
    \begin{equation}
      \rho(\vec{x}) := \left(1-\norm{\vec{x}}^{2}\right)^{m - \frac{\alpha + d}{2}} \sum_{k=0}^{N-1} \rho_k P_{k}^{\left(m - \frac{\alpha + d}{2},\frac{d-2}{2}\right)}\left(2 \norm{\vec{x}}^{2}-1\right)\,,
      \label{eq:ansatz}
    \end{equation}
    with $P_k^{(a, b)}$ the $k$th Jacobi polynomial and $\rho_k$ its corresponding coefficient.
  \end{frame}

  \begin{frame}{Spectral Method Operator}
    \begin{definition}[Power Law Operator $\mathcal{Q}^\beta$]
      The power law operator $\mathcal{Q}^\beta: \functionspace \to \functionspace$ is given by
      $$\mathcal{Q}^\beta[\rho](\vec{x}) := \int \norm{\vec{x}-\vec{y}}^\beta \,\dd\rho(\vec{y}) = \int_{\supp(\rho)} \norm{\vec{x}-\vec{y}}^\beta \rho(\vec{y}) \,\dd\vec{y}\,.$$
    \end{definition}

    Applied to the ansatz given in \eqref{eq:ansatz}, we can evaluate the appearing integrals \alert{explicitly}.
    For the attractive-repulsive interaction kernel $K_{\alpha,\beta}(r) = \frac{r^\alpha}{\alpha} - \frac{r^\beta}{\beta}$, the matrix representation of the operator becomes
    \begin{equation}
      Q_{\alpha, \beta} := \frac{R^{\alpha+d}}{\alpha} Q^\alpha - \frac{R^{\beta+d}}{\beta} Q^\beta\,.
      \label{eq:full-attrep-operator}
    \end{equation}
  \end{frame}

  \begin{frame}{Analysis: The Attractive-Repulsive Operator}
    \begin{figure}[H]
      \centering
      \includegraphics[width=0.9\linewidth]{results/attrep/attractive-repulsive-operators.pdf}
      \caption[Attractive and repulsive operators.]{The attractive and repulsive operators (matrices), the (absolute) matrix values are shown in a $\log_{10}$ colour scale. Due to the choice of basis, the attractive operator is exactly banded.}
      \label{fig:attractive-repulsive}
    \end{figure}
  \end{frame}
  \begin{frame}{Analysis: Comparison with Analytical Solution}
    \begin{columns}[t]
      \begin{column}{0.6\textwidth}
        \begin{figure}[H]
          \centering
          \includegraphics[width=\linewidth, trim={0 6.35cm 0 0}, clip]{results/known-analytic/analytic-solution.pdf}
          \caption[Comparison with analytical solutions and error]{
            The analytical solution $\rho(x)$ compared to the (spectral method) solutions $\rho_N(x)$ of different order.
          }
          \label{fig:analytic-solution}
        \end{figure}
      \end{column}
      \begin{column}{0.4\textwidth}
        \vspace{1.5mm}
        \begin{figure}[H]
          \centering
          \includegraphics[width=\linewidth]{results/known-analytic/simulation-quiver.pdf}
          \caption[Comparison with analytical solutions and error]{Simulation output with $N_p = 150$ particles using the same attractive-repulsive kernel $K_{\alpha,\beta}(r)$.}
        \end{figure}
        % $$K_{\alpha,\beta}(r) = \frac{r^{2.0}}{2.0} - \frac{r^{1.5}}{1.5}\,.$$
      \end{column}
    \end{columns}
  \end{frame}

  \begin{frame}{Original Contribution: General Kernel Spectral Method}
    We express the general kernel as a polynomial, through reprojection from the Jacobi polynomials,
    \begin{equation}
      K_G(r) = \sum_{l=0}^{G-1} g_l r^l \approx K(r)\,,\quad \vec{g} := ({g}_0, ..., {g}_{G-1})^T \in \R^G\,.
      \label{eq:monomial-expansion}
    \end{equation}

    The operator can then be expressed as
    \begin{equation*}
      \mathcal{Q}_G[\hat{\rho}](\hatvec{x})
      = \int_{B_R(\vec{0})} \sum_{l=0}^{G-1} g_l \norm{\hatvec{x}-\hatvec{y}}^l \hat{\rho}(\hatvec{y}) \,\dd\hatvec{y}
      = \sum_{l=0}^{G-1} g_l R^{l+d} \mathcal{Q}^l[\hat{\rho}](\vec{x})\,.
    \end{equation*}
  \end{frame}

  \begin{frame}{Results: Simulation and Solver Comparison}
    \begin{figure}[H]
      \centering
      \includegraphics[width=0.58\linewidth]{results/morse/simulation-solver-comparison.pdf}
      \caption[Comparison of histogram and spectral method solution]{Comparison of simulation output with the $G = 8$ general kernel solver's equilibrium measure $\rho_{12}(r)$ at $R$ given by the simulator.}
      \label{fig:simulation-solver-comparison}
      % for now
    \end{figure}
  \end{frame}

  \begin{frame}{Contributions, Challenges Faced \& Future Work}
    Contributions:
    \begin{itemize}
      \tightlist
      \item Original implementation of simulator (+\glstext{gui}, in C++) and solver (in Julia).
      \item General kernel spectral method + implementation.
      \item Lemma for an initial guess of $R$.
    \end{itemize}

    Challenges:
    \begin{itemize}
      \tightlist
      \item Proving Theorem 4.2, the power law integral of the Jacobi polynomials.
      \item Synchronisation of simulation and solver (implementations).
      \item Correct parameter choices for the general kernel spectral method.
    \end{itemize}

    Future Work:
    \begin{itemize}
      \tightlist
      \item Describing phase-space distributions in a self-propulsion setup.
      \item Modelling (constrained) boundaries and boundary conditions.
    \end{itemize}
  \end{frame}

  \begin{frame}{}
    Thank you.
  \end{frame}

  \begin{frame}[allowframebreaks]
    \frametitle{Bibliography}
    \printbibliography[heading=bibnumbered]
    \printnoidxglossary[type=acronym, title={Acronyms}]
  \end{frame}
\end{document}
